% !PRINTED 407
\weberpage{407}
Sind $\tau', \tau''$ irgend zwei dieser Wurzeln, so folgt aus (3):
\[
\begin{aligned}
\varphi_h\,(\tau')\,\psi_h\,(\tau) - \varphi_h\,(\tau)\,\psi_h\,(\tau') &= 0\\
\varphi_h\,(\tau'')\,\psi_h\,(\tau) - \varphi_h\,(\tau)\,\psi_h\,(\tau'') &= 0,
\end{aligned}
\]
und daraus, da $\psi_h\,(\tau)$ und $\varphi_h\,(\tau)$ nicht zugleich
verschwinden können,
\[
\varphi_h\,(\tau')\,\psi_h\,(\tau'') - \varphi_h\,(\tau'')\,\psi_h\,(\tau') = 0,
\]
d.\,h. es ist, wenn $\tau', \tau''$ irgend zwei der Grössen (7) sind,
\[
\Phi_h\,(\tau', \tau'') = 0.
\]

Daraus folgt, dass die $\nu$ Functionen $\Phi_h\,(t,\tau')$, von einem von
$t$ unabhängigen Factor abgesehen, den nämlichen grössten
gemeinschaftlichen Theiler haben, wie die Functionen $\Phi_h\,(t,\tau)$.
Das Gleiche gilt für die Functionen $\Phi_h\,(t,\tau'')$ u.\,s.\,f., oder
die Gleichungen
\[
R\,(t,\tau) = 0, \quad R\,(t,\tau') = 0, \quad R\,(t,\tau'') = 0, \ldots
\]
haben alle dieselben Wurzeln. Daraus folgt nach (5):
\[
\vartheta = \dfrac{R\,(a,\tau)}{R\,(b,\tau)} = \dfrac{R\,(a,\tau')}{R\,(b,\tau')} = \dfrac{R\,(a,\tau'')}{R\,(b,\tau'')} = \cdots,
\]
und mithin sind nach (6) die $\mu$ Werthe
\[
\xi = \tau, \tau', \tau'', \ldots
\]
die Wurzeln der Gleichung $X = 0$. Andererseits folgt aus
$\Phi_h\,(\tau',\tau) = 0$, $\Phi_h\,(\tau'',\tau) = 0, \ldots$
\[
v_h = \dfrac{\varphi_h\,(\tau)}{\psi_h\,(\tau)} = \dfrac{\varphi_h\,(\tau')}{\psi_h\,(\tau')} = \dfrac{\varphi_h\,(\tau'')}{\psi_h\,(\tau'')} \cdots = \dfrac{1}{\mu}\,\Sigma\,\dfrac{\varphi_h\,(\tau)}{\psi_h\,(\tau)},
\]
und es kann folglich $v_h$ als symmetrische Function der Wurzeln von (6)
rational durch die Coëfficienten dieser Gleichung, d.\,h. rational durch
$\vartheta$ dargestellt werden. Vertauscht man wieder $t$ mit $\tau$, so
erhält man nach (1) und (2) Ausdrücke von der Form

\weq{10}{\begin{aligned}
u &= f\,(\vartheta), \quad u_1 = f_1\,(\vartheta), \ldots, u_{\nu-1} = f_{\nu-1}(\vartheta)\\
\vartheta &= \chi\,(u, u_1, \ldots, u_{\nu-1}),
\end{aligned}}

worin $f, f_1, \ldots, f_{\nu-1}, \chi$ rationale Functionen bedeuten.

\begin{center}
\S.\ 108.\\
\gesperrt{Resolventen mit einem Parameter.}
\end{center}

Wir kehren jetzt zu unseren anfänglichen Voraussetzungen (\S.\,106)
zurück, und bezeichnen mit $u, u_1, \ldots, u_{\nu-1}$ ein System
%CONT>>
